% Options for packages loaded elsewhere
\PassOptionsToPackage{unicode}{hyperref}
\PassOptionsToPackage{hyphens}{url}
%
\documentclass[
]{article}
\usepackage{amsmath,amssymb}
\usepackage{iftex}
\ifPDFTeX
  \usepackage[T1]{fontenc}
  \usepackage[utf8]{inputenc}
  \usepackage{textcomp} % provide euro and other symbols
\else % if luatex or xetex
  \usepackage{fontspec}
  \defaultfontfeatures{Scale=MatchLowercase}
  \defaultfontfeatures[\rmfamily]{Ligatures=TeX,Scale=1}
\fi
\usepackage{lmodern}
\ifPDFTeX\else
  % xetex/luatex font selection
  \setmainfont[]{DejaVu Serif}
\fi
% Use upquote if available, for straight quotes in verbatim environments
\IfFileExists{upquote.sty}{\usepackage{upquote}}{}
\IfFileExists{microtype.sty}{% use microtype if available
  \usepackage[]{microtype}
  \UseMicrotypeSet[protrusion]{basicmath} % disable protrusion for tt fonts
}{}
\makeatletter
\@ifundefined{KOMAClassName}{% if non-KOMA class
  \IfFileExists{parskip.sty}{%
    \usepackage{parskip}
  }{% else
    \setlength{\parindent}{0pt}
    \setlength{\parskip}{6pt plus 2pt minus 1pt}}
}{% if KOMA class
  \KOMAoptions{parskip=half}}
\makeatother
\usepackage{xcolor}
\usepackage[margin=1in]{geometry}
\usepackage{longtable,booktabs,array}
\usepackage{calc} % for calculating minipage widths
% Correct order of tables after \paragraph or \subparagraph
\usepackage{etoolbox}
\makeatletter
\patchcmd\longtable{\par}{\if@noskipsec\mbox{}\fi\par}{}{}
\makeatother
% Allow footnotes in longtable head/foot
\IfFileExists{footnotehyper.sty}{\usepackage{footnotehyper}}{\usepackage{footnote}}
\makesavenoteenv{longtable}
\setlength{\emergencystretch}{3em} % prevent overfull lines
\providecommand{\tightlist}{%
  \setlength{\itemsep}{0pt}\setlength{\parskip}{0pt}}
\setcounter{secnumdepth}{-\maxdimen} % remove section numbering
\usepackage{esint}
\usepackage{mathrsfs}
\emergencystretch=2em
\ifLuaTeX
  \usepackage{selnolig}  % disable illegal ligatures
\fi
\IfFileExists{bookmark.sty}{\usepackage{bookmark}}{\usepackage{hyperref}}
\IfFileExists{xurl.sty}{\usepackage{xurl}}{} % add URL line breaks if available
\urlstyle{same}
\hypersetup{
  hidelinks,
  pdfcreator={LaTeX via pandoc}}

\author{}
\date{}

\begin{document}

\hypertarget{o10-an-unbounded-curvature-rate-exclusion-from-the-fifth-line-test}{%
\section{O10: an unbounded-curvature rate exclusion from the fifth-line
test}\label{o10-an-unbounded-curvature-rate-exclusion-from-the-fifth-line-test}}

2026-09-29. This note continues the existing O10 endpoint analysis. It
does not prove the density-free uniform gap. Unlike a fixed-density
theorem, its hypothesis allows the normalized opposite Weyl energy to
tend to infinity. Its quantitative conclusion is an obstruction to a
definite class of simultaneous endpoint-error/curvature blow-up rates.

\hypertarget{primitive-hypotheses-and-what-is-not-assumed}{%
\subsection{1. Primitive hypotheses and what is not
assumed}\label{primitive-hypotheses-and-what-is-not-assumed}}

All integrals and norms use normalized volume. Write \[
 b_j=\fint|W^-_j|^2,\qquad H_j=(1+b_j)^{1/2},\qquad a_j=|v_j|^2,\qquad q_j=|\nabla v_j|^2.
\] The inherited strong-current and bounded-resolvent results supply,
without a bound on \(b_j\), \[
 |\Phi_j|+|v_j|+|\nabla^*\nabla v_j|\le C,
 \quad \|\nabla\Phi_j\|_4+\|\nabla^2\Phi_j\|_2\le C,
 \quad \|\nabla v_j\|_2\le C.
\] Fix the lower-amplitude cutoffs once and for all, using thresholds
strictly below \(\rho_*\) and \(m_*\). Define \(\eta_j\) to be the
\textbf{sum} of the following measured primitive error norms (and
\(\beta_j^{-1}\)); by the inherited conclusions it tends to zero. This
is a fixed error aggregate, not a freely enlarged sequence with an
adjustable decay rate:

\begin{itemize}
\tightlist
\item
  \(\|a_j-m_*\|_2\), \(\|\rho_j-\rho_*\|_2\), with \(\rho=|\Phi|^2\),
  \(m_*=128/81\), \(\rho_*=4/3\);
\item
  \(\|(\nabla^*\nabla-2)v_j\|_2\), \(\|D^-v_j\|_2\), and the small
  \((1,0)\)-projection \(N_{v_j}\) on a fixed positive-\(\rho\) region;
\item
  \(\|v_j-\tau_j\|_{H^1}+\|v_j-\tau_j\|_4\); the positive-spinor phase
  norm \(\|\chi_\rho p_j\|_2\); and
  \(\|\chi_\rho\chi_a\epsilon_{\eta,j}\|_2\), where
  \(\epsilon_\eta=\eta_{\rm conn}-[\sqrt a/(2\rho^{3/2})]\theta_1\) on
  the fixed positive-amplitude region. These norms are invariant under
  the compatible adapted-frame phase changes; the cutoffs make their
  global definition smooth.
\item
  the actual determinant-curvature fit \(\| -iF_A^-+(27/32)Q_v\|_2\),
  the traceless Ricci error, and the positive-block curvature errors
  used in the inherited Cartan formulas.
\end{itemize}

The second bullet is inherited from the actual first-order system and
the \(H^1\) closeness of \(v\) to \(\tau\). The first-connection error
in the third bullet is obtained algebraically from the actual cubic
formula and \(v-\tau\); it is not a derivative of the negative-spinor
connection. Those approximation errors are explicitly included in the
fixed sum above. An equivalent maximum of this same finite list changes
only absolute constants.

\textbf{The definition does not include} \(da_j\), the negative-spinor
phase \(j_{v_j}\), the first prolongation error \(\epsilon_b\),
\(D_5c_j\), or any claimed convergence of \(q_j\). Those are conclusions
below. No rate formula expressing \(\eta_j\) in terms of the original
scalar defect \(\kappa_j\) is established in this note.

The derivative-interchange estimate gives, with an absolute inherited
constant, \[
 \|\nabla^2v_j\|_2+\|\nabla v_j\|_4^2\le C H_j. \tag{1}
\] This is an estimate on each smooth member, not a uniform bound along
the sequence.

\hypertarget{result}{%
\subsection{2. Result}\label{result}}

Subject to the inherited first-order identities and the fifth-line
Cartan/symbol identities, an actual endpoint sequence cannot satisfy \[
 \boxed{\eta_j(1+b_j)^{3/2}=\eta_j H_j^3\longrightarrow0.} \tag{2}
\] Thus a hypothetical unrestricted O10 escape must have, for this
primitive error aggregate, \[
 \liminf_j\eta_j(1+b_j)^{3/2}>0. \tag{3}
\] This is a sequence restriction, not a constant independent of density
in the original O10 inequality. Its exponent is sufficient for the
proof; no claim of optimality is made.

\hypertarget{radial-derivative-and-the-surviving-connection-component}{%
\subsection{3. Radial derivative and the surviving connection
component}\label{radial-derivative-and-the-surviving-connection-component}}

The exact scalar identity is \[
 \Delta a=4a-2q+e_a,\qquad \|e_a\|_2\le C\eta.
\] Using (1), \(\|\Delta a\|_2\le C H\). Testing with \(a-m_*\)
therefore gives \[
 \boxed{\|da\|_2^2\le C\eta H.} \tag{4}
\] The exact mass-phase identity on the fixed positive-\(\rho\) region
gives \[
 |j_v+\tfrac12J_+da|\le C|N_v|.
\] The exceptional sets where \(\rho\) or \(a\) lies below a fixed
positive cutoff have normalized measure at most \(C\eta^2\). By (1), \[
 \fint_{\mathrm{bad}}|\nabla v|^2
 \le \mu(\mathrm{bad})^{1/2}\|\nabla v\|_4^2
 \le C\eta H. \tag{5}
\] Consequently \[
 \|j_v\|_2\le C(\eta+\eta^{1/2}H^{1/2}).
\] The same mass-phase-Dirac reduction as in the finite-density proof
now gives, on the fixed positive-amplitude regions, \[
 b=c_0\theta_2+\epsilon_b,
 \qquad \|\epsilon_b\|_2\le C(\eta+\eta^{1/2}H^{1/2}). \tag{6}
\] Here \(b\) is the off-diagonal connection one-form, not the scalar
Weyl budget. Formula (6) follows from the small \((1,0)\) derivative,
the first spinor components \((da/2+ij_v)/\sqrt a\), and the Dirac
trace. It does not require differentiating \(N_v\), or a Nijenhuis
bracket estimate.

Define the global cut-off fifth-line section exactly as in the supplied
proof: \[
 c=\chi_\rho(\rho)\chi_a(a)
 \left(-\frac{\Phi\wedge t(v)}{2\sqrt\rho\,a^{3/2}}\right)\in\Gamma(L^5).
\] It satisfies \[
 \|c\|_2\le C,\qquad \|c\|_4\le C H^{1/2},
 \qquad \|D_5c\|_2\le C H. \tag{7}
\] The bounded denominators in this assertion are imposed by the fixed
cutoffs; there is no division at an uncontrolled zero.

\hypertarget{a-covariance-estimate-prevents-loss-of-the-fifth-line-mass-into-its-tail}{%
\subsection{4. A covariance estimate prevents loss of the fifth-line
mass into its
tail}\label{a-covariance-estimate-prevents-loss-of-the-fifth-line-mass-into-its-tail}}

Set \(q_*=2m_*\) and \(T(q)=q/(1+q)\). Then \[
 |dT(q)|\le C|\nabla^2v|.
\] Testing the scalar identity with \(T(q)-T(q_*)\), and using (1), (4),
gives \[
 \boxed{\fint\frac{(q-q_*)^2}{1+q}
 \le C\bigl(\eta+\eta^{1/2}H^{3/2}\bigr).} \tag{8}
\] The term with \(da\) is bounded by
\(\|da\|_2\|dT(q)\|_2\le C\eta^{1/2}H^{3/2}\). This is the largest power
of \(H\) in the argument.

Since \(\fint q\le C\), Cauchy--Schwarz gives \[
 \|q-q_*\|_1
 \le C\left(\eta+\eta^{1/2}H^{3/2}\right)^{1/2}\longrightarrow0 \tag{9}
\] under (2). This step needs neither a uniform \(L^2\) bound on \(q\)
nor a fixed bound on \(b\).

The surviving-component norm identity, (4)--(6), and \[
 \fint|a-m_*||c|^2\le\|a-m_*\|_2\|c\|_4^2\le C\eta H
\] use the exact orthogonal-projection identity
\(|b|^2=2|c_0|^2+|\epsilon_b|^2\). On the positive-amplitude region, \[
 q-2a|c_0|^2=|da|^2/(4a)+a|q_{\mathrm{phase}}|^2+a|\epsilon_b|^2\ge0.
\] Equations (4)--(6) and the cutoff-region estimate give
\(\|q-2m_*|c|^2\|_1\le C\eta H\). They therefore imply the explicit
estimate \[
 \||c|^2-1\|_1\le C\left[(\eta+\eta^{1/2}H^{3/2})^{1/2}+\eta H\right].
\] In particular, \[
 \boxed{\||c|^2-1\|_1\longrightarrow0.} \tag{10}
\] This is more useful here than the old tail estimate
\(\int_{|c|>M}|c|^2\le C_b/M^2\), whose constant blows up. For every
fixed \(M>1\), (10) gives \[
 \fint_{|c|>M}|c|^2
 \le (1-M^{-2})^{-1}\||c|^2-1\|_1\longrightarrow0. \tag{11}
\] Therefore the positive line resolvent may retain a \textbf{fixed}
truncation level, even though \(b_j\to\infty\).

\hypertarget{the-weak-cr-errors-still-vanish-in-this-regime}{%
\subsection{5. The weak CR errors still vanish in this
regime}\label{the-weak-cr-errors-still-vanish-in-this-regime}}

Use the same exact Cartan equation and symbol identity as the supplied
fifth-line proof: \[
 (D_5c_0\wedge\theta_2)^+
 =-(d_{3A}\epsilon_b)^++G,
 \qquad \mathcal L^*\mathcal D\mathcal L=I-iJ_-.
\] Let \(w\in\Gamma(L^5)\) satisfy fixed bounds on
\(\|w\|_\infty+\|D_5w\|_4+\|D_5^2w\|_2\). Tracking the factors rather
than replacing them by \(C_b\) gives:

\begin{longtable}[]{@{}
  >{\raggedright\arraybackslash}p{(\columnwidth - 2\tabcolsep) * \real{0.5000}}
  >{\raggedright\arraybackslash}p{(\columnwidth - 2\tabcolsep) * \real{0.5000}}@{}}
\toprule\noalign{}
\begin{minipage}[b]{\linewidth}\raggedright
Error
\end{minipage} & \begin{minipage}[b]{\linewidth}\raggedright
Bound up to a test-dependent absolute constant
\end{minipage} \\
\midrule\noalign{}
\endhead
\bottomrule\noalign{}
\endlastfoot
Derivative moved off \(\epsilon_b\), hitting \(D_5w\) &
\(\eta+\eta^{1/2}H^{1/2}\) \\
Derivative hitting the adapted tensor map &
\((\eta+\eta^{1/2}H^{1/2})(1+H^{1/2})\) \\
Negative phase times \(c\), paired with \(D_5w\) &
\((\eta+\eta^{1/2}H^{1/2})H^{1/2}\) \\
\(c\,d\chi\) on a low-amplitude transition region &
\(\mu(\mathrm{bad})^{1/4}\|c\|_4\|d\chi\|_4\le C\eta^{1/2}H\) \\
Replacing \(J_-\) by \(\bar J=(a/m_*)J_-\) &
\(\|a-m_*\|_4\|D_5c\|_2\le C\eta^{1/2}H\) \\
Bochner drift \(\delta\bar\Omega\,c\,D_5w\) &
\(C(\eta+\eta^{1/2}H^{1/2})H^{1/2}\) \\
\end{longtable}

Here \(\|a-m_*\|_4\le C\eta^{1/2}\), because \(v\) is uniformly bounded.
The positive-spinor derivative factors have fixed \(L^4\) bounds. The
potential error is \(O(\eta)\), paired with \(c\in L^2\) and bounded
\(w\).

All displayed errors tend to zero under (2). Hence the weak conclusion
survives: \[
 \Re\fint\langle c,(\Delta_{5A}+10/3)w\rangle\longrightarrow0. \tag{12}
\] This calculation does \textbf{not} make (12) valid for an arbitrary
unrestricted sequence. Without a relation such as (2), the factor
\(\epsilon_b\nabla v\,D_5w\) remains uncontrolled.

\hypertarget{fixed-level-resolvent-contradiction}{%
\subsection{6. Fixed-level resolvent
contradiction}\label{fixed-level-resolvent-contradiction}}

Fix \(M>2\) and solve \[
 (\Delta_{5A}+10/3)w=T_M(c).
\] The right side is uniformly bounded by \(2M\). The maximum principle,
the original Ricci-energy bound, and the controlled curvature \(5F_A\)
give all the fixed test bounds required above, with no
\(b_j\)-dependence. Thus (12) gives \[
 \Re\fint\langle c,T_M(c)\rangle\longrightarrow0.
\] But (10)--(11) imply the left side tends to 1. This contradiction
proves (2) is impossible.

\hypertarget{exact-scope-and-a-rejected-shortcut}{%
\subsection{7. Exact scope and a rejected
shortcut}\label{exact-scope-and-a-rejected-shortcut}}

This calculation isolates two distinct difficulties in the fully
unrestricted problem:

\begin{enumerate}
\def\labelenumi{\arabic{enumi}.}
\tightlist
\item
  radial/derivative concentration prevents \(|c|^2\) from having a
  nonzero uniformly integrable limit;
\item
  critical weak CR products couple a small first-order error to a
  growing fourth-power derivative budget.
\end{enumerate}

The covariance estimate controls the first and the explicit product
estimates control the second under the same sufficient condition (2).
This is a quantitative extension beyond a fixed finite density, but it
is not a density-free theorem until the actual equations force (2), or
supply a different contradiction in its complement.

No explicit modulus \(\eta\le\omega(\kappa,\beta^{-1})\) with a power
rate is claimed here. The inherited strong-current argument supplies
smallness; extracting a modulus requires tracing its truncation and
resolvent choices. Even an explicit modulus alone would not bound \(b\)
or force (2).

A possible shorter route through the positive-spinor curvature was
checked and fails algebraically. In the ideal adapted configuration
\(\eta_{\rm conn}=\kappa\theta_1\), \(b=c\theta_2\), the exact Cartan
formula is \[
 d\theta_1=-i(\mathfrak a+p-q)\wedge\theta_1+\kappa\Xi_+,
 \qquad \Xi_+=e^{13}-e^{24}+i(e^{14}+e^{23}).
\] The two real components of \(c\) cancel exactly. Consequently \[
 (F^{S^+})_{21}
 =d\kappa\wedge\theta_1+i\kappa(q-3p)\wedge\theta_1+\kappa^2\Xi_+.
\] The controlled positive curvature does not directly bound or
eliminate the surviving \(c\). This finite algebra was independently
checked using the Clifford-to-Cartan script; it prevents mistaking a
missing negative-curvature estimate for an immediately available
positive-curvature one.

\hypertarget{exact-obstruction-to-the-proposed-amplitude-weighted-shortcut}{%
\subsection{8. Exact obstruction to the proposed amplitude-weighted
shortcut}\label{exact-obstruction-to-the-proposed-amplitude-weighted-shortcut}}

The following is an exact local calculation in the zero-error
first-order system, retaining the radial derivatives rather than
assuming them small. It explains why the proposed replacement of \(c\)
by \(m^{-3/2}c\) does not yet remove the density dependence. It is not a
new limiting argument.

Write \(m=|v|^2>0\), \(m_i=e_i(m)\), and use the full directional fields
\(e_1-ie_2,e_3-ie_4\); changing both to half-sized fields does not
change the final real-component formulas. In an adapted frame
\(v=\sqrt m(1,0)\), the equations \(N_v=0\), \(D^-v=0\) give \[
 q_{\rm ph}=\frac1{2m}(m_2,-m_1,m_4,-m_3)
 =-\frac{J_+dm}{2m},
\] \[
 \boxed{b=d\,\theta_1+c\,\theta_2,
 \qquad d=-\frac{m_3+i m_4}{2m}.} \tag{14}
\] In particular, one may not begin this calculation with
\(b=c\theta_2\): doing so already discards the plane-two radial
derivative.

The additional \(d\)-terms in the exact Clifford-to-Cartan conversion
are \[
 d\theta_1=-i(\mathfrak a+p-q_{\rm ph})\wedge\theta_1
             +\kappa\Xi_+-d\,\overline{\Xi_+},
\] \[
 d\theta_2=i(2\mathfrak a-p-q_{\rm ph})\wedge\theta_2
             -\bar c\Xi_- -2i\bar d\,e^{12}.
 \tag{15}
\] Here \(\overline{\Xi_+}=\theta_1\wedge\theta_2\). The added
coefficients were checked separately on the real and imaginary basis
directions of \(d\).

Set \(\rho=\rho_*\), \(p=0\), and \(\kappa=\sqrt m/(2\rho_*^{3/2})\), to
isolate the proposed mass mechanism. The positive off-diagonal curvature
is \[
 (F^{S^+})_{21}
 =(d\kappa+i\kappa q_{\rm ph})\wedge\theta_1
  +\kappa^2\Xi_+-\kappa d\,\overline{\Xi_+}.
\] Since \[
 dm-iJ_+dm=(m_1+i m_2)\theta_1+(m_3+i m_4)\theta_2,
\] its first term equals \(-\kappa(m_3+i m_4)\overline{\Xi_+}/(2m)\).
The last term has the opposite coefficient by (14). Thus \[
 \boxed{(F^{S^+})_{21}=\kappa^2\Xi_+.} \tag{16}
\] \textbf{The radial term cancels too.} The controlled positive
curvature therefore does not first force \((dm)_{\{3,4\}}\) to be small.
Such a prior conclusion cannot be used to replace \(J_+dm\) by \(J_-dm\)
and absorb the entire phase into \(m^{-3/2}c\).

The independent mass problem is equally explicit. Put
\(E_1=m_1^2+m_2^2\), \(E_2=m_3^2+m_4^2\). The full first-jet norm in
(14) is \[
 \boxed{|\nabla v|^2=\frac{E_1+2E_2}{2m}+2m|c|^2.} \tag{17}
\] In the exact eigenvalue equation \(\nabla^*\nabla v=2v\), testing
\(\Delta m=4m-2|\nabla v|^2\) with \(m\) consequently gives \[
 \boxed{2\int E_1+3\int E_2+4\int m^2|c|^2
             =4\int m^2.} \tag{18}
\] A lower bound for \(\int m|\nabla v|^2\) alone therefore does not
produce a positive lower bound for \(\int m^2|c|^2\). The radial terms
can pay the entire right side of (18). This is exactly the term that the
finite-density mass-gradient estimate removes; a positive polynomial
weight does not silently remove it.

More generally, the same formal exact system gives for positive integer
\(p\) \[
 (p+1)\int m^{p-1}E_1+(p+2)\int m^{p-1}E_2
       +4\int m^{p+1}|c|^2=4\int m^{p+1}.
\] These identities do not by themselves contradict an amplitude
concentration defect: for example their limiting radial moment budget
with no \(c\)-mass is algebraically compatible with the positive measure
\(d\mu_1(t)=4t\,dt\) on \([0,m_*]\), \(\mu_2=0\). This is only an
illustration of what the moment identities fail to exclude, \textbf{not}
a construction of actual metrics or fields.

Accordingly the weighted shortcut is not closed. Its exact unresolved
terms are the plane-two radial coefficient in the Cartan equation and
the radial concentration contribution in the mass identity. Neither is
an omitted harmless cutoff error.

\end{document}
